\documentclass[11pt]{amsart}
\usepackage[margin=1.08in]{geometry}
\usepackage{amsmath,amssymb,amsthm,mathtools,microtype}
\usepackage{hyperref}
\usepackage{xurl}
\hypersetup{colorlinks=true,linkcolor=blue,citecolor=blue,urlcolor=blue,
 pdftitle={Global regularity for critical stochastic SQG with finite-dimensional nonlinear noise},
 pdfauthor={Kaiwen Jin}}
\newtheorem{theorem}{Theorem}[section]
\newtheorem{proposition}[theorem]{Proposition}
\newtheorem{lemma}[theorem]{Lemma}
\theoremstyle{remark}
\newtheorem{remark}[theorem]{Remark}
\numberwithin{equation}{section}
\newcommand{\T}{\mathbb T}
\newcommand{\R}{\mathbb R}
\newcommand{\PP}{\mathbb P}
\newcommand{\EE}{\mathbb E}
\newcommand{\dd}{\,\mathrm d}
\newcommand{\norm}[1]{\lVert#1\rVert}
\newcommand{\fint}{\mathchoice{\displaystyle\int\!\!\!\!-}{\int\!\!\!-}{\int\!-}{\int\!-}}
\title[Critical stochastic SQG]{Global regularity for critical stochastic SQG
with finite-dimensional nonlinear noise}
\author{Kaiwen Jin}
\thanks{Contact: \href{mailto:yui.ui.siki@gmail.com}{yui.ui.siki@gmail.com}.}
\date{October 2026}
\keywords{Surface quasi-geostrophic equation, critical dissipation, nonlinear multiplicative noise, global regularity}
\begin{document}
\begin{abstract}
We prove almost-sure global regularity and pathwise uniqueness for the critical
surface quasi-geostrophic equation on the two-dimensional torus with a finite
collection of smooth nonlinear state-dependent It\^o noise coefficients.
For every deterministic smooth datum, globally scalar-Lipschitz coefficients
with bounded positive-order derivatives give a global adapted solution
continuous in $H^3$ on every finite time interval, with the integral equation
holding in $H^1$. There is no smallness, mean-zero or affine-noise condition.
The proof integrates finite powers of spatial increments. A dissipation bound
proportional to the power permits a spatial exponent independent of that power,
so a sufficiently large finite power controls a positive H\"older norm.
A fractional stochastic truncation argument bounds the amplitude, and a tame
$H^3$ estimate then prevents finite-time blowup.
\end{abstract}
\maketitle

\section{Introduction}\label{sec:intro}

The surface quasi-geostrophic equation is a nonlocal active scalar model in
geophysical fluid dynamics. Its velocity is recovered from the temperature by
Riesz transforms. With order-one dissipation, transport and diffusion act at the
same differential order, and an energy estimate alone does not control smooth
solutions. Deterministic critical SQG nevertheless has global regularity;
Constantin, Tarfulea and Vicol \cite{CTV}, for example, obtain it by estimating
spatial increments and propagating a small positive H\"older exponent from an
amplitude bound.

We consider the It\^o equation
\begin{equation}\label{eq:sqg}
 \dd\theta+\big(R^\perp\theta\cdot\nabla\theta+\Lambda\theta\big)\dd t
       =\sum_{k=1}^K g_k(x,\theta)\dd W^k_t,
 \qquad \theta(0)=\theta_0,
\end{equation}
on $\T^2$, where $\Lambda=(-\Delta)^{1/2}$ and
$R^\perp=(-R_2,R_1)$, with zero Fourier multipliers on constants.
The $W^k$ are independent standard Brownian motions on a prescribed filtered
probability space satisfying the usual conditions. The datum $\theta_0$ is any
deterministic smooth real periodic function. We assume that the finite family
$g_k\in C^\infty(\T^2\times\R;\R)$ is globally Lipschitz in its scalar
argument, uniformly in $x$, that $g_k(\cdot,0)$ is bounded, and that
\begin{equation}\label{eq:coefficients}
 \sup_{x,r}|\partial_x^\beta\partial_r^j g_k(x,r)|<\infty
 \qquad\text{whenever }|\beta|+j\geq1.
\end{equation}
Each derivative may have its own bound. In particular the noise can depend
nonlinearly on the temperature and smoothly on space.

Local pathwise theory for stochastic fractional SQG is available from
R\"ockner, Zhu and Zhu \cite{RZZ}. Global results at subcritical dissipation
are discussed in Agresti and Veraar \cite[Theorem 8.22]{AV}; their
Remark 8.23 identifies the critical endpoint as a remaining question for the
general coefficient class there. More recently, Chen and Zhang
\cite[Section 6]{CZ} proved critical stochastic well-posedness on $\R^2$ with
additive Brownian forcing. A direct replacement of that additive forcing by
$g(x,\theta)$ would require an a priori spatial regularity bound for the
composed coefficient, which the energy estimate does not provide.
The finite smooth nonlinear class \eqref{eq:coefficients} gives a concrete
critical problem beyond this additive setting. It is the class specified in
Problem 478 of the AIM catalogue \cite{AIM478}.

\begin{theorem}\label{thm:main}
For every datum and finite coefficient family just described, equation
\eqref{eq:sqg} has a pathwise unique global adapted solution on the prescribed
stochastic basis. Almost surely, simultaneously for every finite $T$,
\[
 \theta\in C([0,T];H^3(\T^2)),
\]
and, for all $t\in[0,T]$, the integral identity
\[
 \theta(t)=\theta_0-
 \int_0^t\big(R^\perp\theta(s)\cdot\nabla\theta(s)+\Lambda\theta(s)\big)\dd s
 +\sum_{k=1}^K\int_0^t g_k(\cdot,\theta(s))\dd W^k_s
\]
holds in $H^1(\T^2)$. No smallness or condition on the spatial mean is imposed.
\end{theorem}

The theorem gives a negative answer to finite-time blowup with positive
probability in the full smooth finite-noise class of \cite{AIM478}.
In particular, no affine or additive restriction on $g_k$ is introduced.
The conclusion concerns this stated class; the broader rough or
infinite-dimensional coefficient classes in \cite{AV} are not asserted here.

The deterministic increment estimates in \cite{CTV} suggest propagating a
weighted pointwise supremum. For It\^o paths, their direct time-differentiation
argument is unavailable. We instead integrate a finite power of the weighted
increments over space and separation. The scalar convexity remainder retains
a factor proportional to that power. It can therefore absorb the transport of
the weight with a spatial exponent chosen independently of the power. Taking
the power large, but finite, then gives a positive H\"older exponent by a
fractional Sobolev estimate. The stochastic trace and bracket are controlled by
the same integrated quantity.

Before fixing an amplitude level, a fractional stochastic truncation argument
shows that the amplitude is finite up to the possible maximal lifetime. This
uses the martingale first-failure method of Hsu, Wang and Wang \cite{HWW},
with the order-one dissipation and the periodic $H^{1/2}\to L^4$ inequality
included explicitly. The positive H\"older bound then makes transport
absorbable in an $H^3$ energy estimate. A tame chain rule controls the nonlinear
noise using the amplitude, and logarithmic energy bounds exclude a finite
maximal lifetime.

Section~\ref{sec:local} gives the maximal local solution and the martingale
bounds. Section~\ref{sec:amplitude} proves the amplitude estimate.
Section~\ref{sec:regularity} establishes the integrated increment estimate
and completes the continuation argument. We use
$\T^2=(\R/(2\pi\mathbb Z))^2$ with ordinary Lebesgue measure and
$\widehat v(m)=(2\pi)^{-2}\int_{\T^2}v(x)e^{-im\cdot x}\dd x$.
The symbols of $\Lambda$ and $R_j$ are $|m|$ and $im_j/|m|$ for $m\ne0$,
and zero for $m=0$.

\section{Local solutions and stochastic estimates}\label{sec:local}

We use real inhomogeneous Sobolev spaces on the torus. Their norms are
\[
 \norm{v}_{H^s}^2=(2\pi)^2\sum_{m\in\mathbb Z^2}
 (1+|m|^2)^s|\widehat v(m)|^2.
\]
The norm in $C^\gamma$, $0<\gamma<1$, includes the supremum and the
$\gamma$-H\"older seminorm. All constants may depend on the fixed coefficient
family and the torus. Other dependencies are indicated locally. Constants in
estimates at a local stopping time will be independent of its level unless
explicitly stated otherwise.

\subsection{The maximal local solution}

The constant mode must remain in the solution space because the noise need not
preserve the mean. The abstract local theorem of R\"ockner, Zhu and Zhu
\cite[Theorem 3.2 and Definition 3.1]{RZZ} accommodates this directly.

\begin{proposition}\label{prop:local}
Under the hypotheses of Theorem~\ref{thm:main}, there is a pathwise unique
maximal local solution $\theta$ on the prescribed stochastic basis, with lifetime
$\tau_*>0$ almost surely. It satisfies
\[
 \theta\in C([0,\tau_*);H^3),\qquad
 \theta\in L^2_{\mathrm{loc}}([0,\tau_*);H^{7/2})
 \quad\hbox{almost surely}.
\]
There are increasing local norm stopping times $\tau_n\uparrow\tau_*$,
with deterministic $H^3$ bounds at each fixed level. If $\tau_*<\infty$,
then $\sup_{t<\tau_*}\norm{\theta(t)}_{H^3}=\infty$.
\end{proposition}

\begin{proof}
Let $P_0$ denote orthogonal projection onto constants in $L^2(\T^2)$ and set
\begin{equation}\label{eq:localrealization}
 A=-\Delta+P_0,\qquad
 f(v)=R^\perp v\cdot\nabla v-P_0v,\qquad
 G(v)e_k=g_k(\cdot,v).
\end{equation}
The operator $A$ is positive self-adjoint on full $L^2$, has domain $H^2$
and compact inverse, and $A^{1/2}=\Lambda+P_0$. Its spectral Sobolev norms
are equivalent to the inhomogeneous norms, including the constant mode.
Thus the source equation $\dd\theta+(f(\theta)+A^{1/2}\theta)\dd t
=G(\theta)\dd W$ is precisely \eqref{eq:sqg}. In the notation of that
theorem the regularity index is $s_0=3$ and the dissipation exponent is $1/2$.

We verify the drift and noise hypotheses. The periodic product, commutator and
Sobolev inequalities in \cite[Lemmas 5.1--5.3 and Section 5.2]{RZZ} give, for
$s\in[3,4]$ and $u=R^\perp v$,
\begin{align}
 |\langle\Lambda^s(u\cdot\nabla v),\Lambda^sv\rangle|
 &\leq C\norm{v}_{H^s}
 \big(\norm{\nabla u}_4\norm{\Lambda^sv}_4
       +\norm{\Lambda^su}_4\norm{\nabla v}_4\big)\notag\\
 &\leq C\norm{v}_{H^{3/2}}\norm{v}_{H^s}\norm{v}_{H^{s+1/2}}\notag\\
 &\leq\varepsilon\norm{v}_{H^{s+1/2}}^2
   +C_\varepsilon\norm{v}_{H^3}^2\norm{v}_{H^s}^2.
 \label{eq:localcoercive}
\end{align}
The first line uses transport cancellation and commutator exponents all equal
to $4$; $H^{1/2}\hookrightarrow L^4$ applies to the differentiated,
zero-mean factors. The mean of $u\cdot\nabla v$ is zero, so its exact
$A$-spectral energy equals the homogeneous pairing, while the additional drift
in \eqref{eq:localrealization} contributes $-\norm{P_0v}_2^2$.
This proves the source coercivity condition in its explicitly stated equivalent
form. For $s=2,3$, the same product estimate gives
\[
 \norm{f(v)}_{H^s}\leq C(1+\norm{v}_{H^3})\norm{v}_{H^{s+3/2}}.
\]
These bounds also give the required continuous drift maps and smooth-test
pairings. For $q=v_1-v_2$, divergence-free transport cancels and leaves
$\langle q,R^\perp q\cdot\nabla v_1\rangle$ and the bounded mean term.
Consequently the local $L^2$ monotonicity condition follows from
$\norm{\nabla v_1}_\infty\leq C\norm{v_1}_{H^3}$ and the Riesz $L^2$ bound.

The coefficient assumptions imply
\begin{equation}\label{eq:localcomposition}
 \sum_k\norm{g_k(\cdot,v)}_{H^s}^2
 \leq C(\norm{v}_{H^3})(1+\norm{v}_{H^s}^2),\qquad 3\leq s\leq4,
\end{equation}
with a locally bounded coefficient, as well as the global $L^2$ Lipschitz
bound for $G$. Here is a verification of \eqref{eq:localcomposition} allowing
full spatial and scalar dependence. Write $s=3+\sigma$, $0\leq\sigma\leq1$.
Every third spatial derivative of $g(x,v(x))$ has the form
$g_r(x,v)D^3v+F$, where $F$ is a sum of bounded-coefficient products of
$DvD^2v$, $(Dv)^3$, $D^2v$, $(Dv)^2$, $Dv$ and a bounded function.
On an $H^3$ ball, $Dv\in L^\infty$, $D^2v\in L^4$ and $D^3v\in L^2$.
Differentiating $F$ once, including the product $(D^2v)^2$, therefore shows
$\norm{F}_{H^1}\leq C(\norm{v}_{H^3})$.
Also $a(x)=g_r(x,v(x))$ has $H^3$ norm bounded on that ball by the integer
chain rule. Fourier convolution gives, uniformly for $0\leq\sigma\leq1$,
\[
 \norm{az}_{H^\sigma}\leq
 C\Big(\sum_m(1+|m|)^\sigma|\widehat a(m)|\Big)\norm{z}_{H^\sigma}
 \leq C\norm{a}_{H^3}\norm{z}_{H^\sigma}.
\]
The last sum is controlled by Cauchy--Schwarz since
$\sum_m(1+|m|)^{-6+2\sigma}<\infty$ uniformly in this interval.
This proves \eqref{eq:localcomposition}; lower derivatives and the $L^2$ term
follow from the same chain rule and linear growth. Finiteness of the family
removes any summability issue. The $L^2$ Lipschitz bound follows directly from
scalar Lipschitz continuity.

All hypotheses of the cited abstract theorem now hold on the original basis.
The datum is deterministic and has finite $H^3$ norm. The theorem and its
maximal-solution definition yield the asserted lifetime and pathwise uniqueness.
Intersecting its local exhaustion with the first $H^3$ level-$n$ hit, if
necessary, gives the indicated norm stops without changing the lifetime.
\end{proof}

The regularity in Proposition~\ref{prop:local} justifies the calculations below
at each fixed local stop. In particular $H^3(\T^2)\hookrightarrow C^{1,\beta}$
for every $0<\beta<1$, for both $\theta$ and $u$. This follows directly from
Cauchy--Schwarz in the Fourier sum with weight
$(1+|m|^2)^{-3}|m|^{2+2\beta}$. No higher local smoothness is required.

\subsection{Martingale bounds}

We use the following two forms of the continuous martingale estimate.

\begin{lemma}\label{lem:martingale}
If $M$ is a continuous local martingale, $a,b>0$ and $S<\infty$, then
\begin{equation}\label{eq:oscillation}
 \PP\left\{\sup_{0\leq s\leq t\leq S}(M_t-M_s)>a,
                  [M]_S\leq b\right\}
 \leq C\exp(-ca^2/b)
\end{equation}
for universal positive constants. Suppose nonnegative continuous semimartingales
$X_n$ have a common deterministic initial value and, before prescribed stops,
\begin{equation}\label{eq:genericlog}
 \dd X_n\leq C(1+X_n)\dd t+\dd M_n,\qquad
 \dd[M_n]_t\leq C(1+X_n)^2\dd t,
\end{equation}
with $C$ independent of $n$. On every fixed horizon, the suprema of the stopped
$\log(1+X_n)$ have a common tail tending to zero, and all finite moments of
the stopped $\sup X_n$ are bounded uniformly in $n$.
\end{lemma}

\begin{proof}
Estimate \eqref{eq:oscillation} is the continuous local martingale oscillation
bound of Hsu, Wang and Wang \cite[Lemma 3.1, arXiv version 3]{HWW}; only unspecified universal
constants are needed here. Applying It\^o's formula to \eqref{eq:genericlog}
gives, on $[0,S]$,
\[
 \log(1+X_n(t))\leq\log(1+X_n(0))+CS+L_n(t),\qquad [L_n]_S\leq CS.
\]
The second derivative term has nonpositive sign. Formula
\eqref{eq:oscillation}, with the initial time included, gives a Gaussian
upper tail for $\sup L_n$. Its exponential moments of every finite order
are finite, proving both assertions. A zero bracket gives the same statements
directly. Auxiliary martingale stops can first be inserted and then removed.
\end{proof}

We apply finite-dimensional or scalar It\^o calculus before taking spatial
limits; see \cite[Theorem 2.4]{BVNVW}. The finite family and the local $H^3$
bounds ensure the required integrability. Stochastic integrals commute with
the spatial integrals below by approximation with simple functions and the
stopped $L^2$ isometry. For an indicator before a stopping time one can use
the predictable version $1_{t\leq\tau_n}$, which gives the same Brownian
integral as $1_{t<\tau_n}$.

\section{An amplitude bound before the maximal lifetime}\label{sec:amplitude}

The coefficient assumptions give
\begin{equation}\label{eq:growth}
 \sum_k|g_k(x,r)|^2\leq C(1+r^2).
\end{equation}
We first prove an amplitude estimate using only this growth bound and the
divergence-free property of the drift. In particular, its constants do not
involve an assumed amplitude bound.

\begin{proposition}\label{prop:amplitude}
For $0<\varepsilon<S<\infty$ and every finite $q\geq1$,
\begin{equation}\label{eq:amplitudemoments}
 \sup_n\EE\left[\sup_{\varepsilon\leq t\leq S,\ t<\tau_n}
             \norm{\theta(t)}_\infty^q\right]<\infty,
\end{equation}
where an empty supremum is zero. Almost surely, for every finite $S$,
\begin{equation}\label{eq:amplitudeaslifetime}
 \sup_{0\leq t<S\wedge\tau_*}\norm{\theta(t)}_\infty<\infty.
\end{equation}
\end{proposition}

\begin{proof}
The $L^2$ It\^o identity, transport cancellation and \eqref{eq:growth} imply
\begin{equation}\label{eq:L2moments}
 \sup_n\EE\sup_{t\leq S\wedge\tau_n}
       (1+\norm{\theta(t)}_2^2)^q<\infty
 \quad(S<\infty, q<\infty).
\end{equation}
Indeed the drift of $X=\norm{\theta}_2^2$ is at most $C(1+X)$ after dropping
$2\norm{\Lambda^{1/2}\theta}_2^2$, and the energy martingale has bracket
at most $C(1+X)^2\dd t$. Lemma~\ref{lem:martingale} applies.

Fix $n$. Set $\widehat\theta_n(t)=\theta(t)$ for $t<\tau_n$ and
$\widehat\theta_n(t)=0$ afterwards. For $a\geq0$ the energy of
$z=(\widehat\theta_n-a)_+$ has a downward jump at the killing time.
Before that time its truncation inequality is
\begin{equation}\label{eq:truncation}
 \dd\norm{z}_2^2+2\norm{\Lambda^{1/2}z}_2^2\dd t
 \leq C\int_{\theta>a}(1+\theta^2)\dd x\dd t
       +2\sum_j\left(\int z g_j(x,\theta)\dd x\right)\dd W^j.
\end{equation}
Every integrand is understood to vanish after killing. The transport term is
$\int u\cdot\nabla(z^2)/2=0$. For the fractional term, integrate
\[
 (r-s)\big((r-a)_+-(s-a)_+\big)
       \geq\big|(r-a)_+-(s-a)_+\big|^2
\]
against the symmetric positive kernel of $\Lambda$.

To justify the nonsmooth truncation, take smooth convex functions $\phi_\eta$
vanishing on $(-\infty,a]$, with $0\leq\phi_\eta''\leq2$ and
$\phi_\eta'\to2(r-a)_+$ uniformly on bounded scalar intervals. Pointwise
It\^o's formula, followed by spatial integration, applies at a fixed $H^3$
stop. The fractional bilinear forms converge by the local Lipschitz bound,
the trace terms are dominated by $C(1+\theta^2)$, and the stochastic terms
converge by the stopped $L^2$ isometry. This proves \eqref{eq:truncation}.
Its integrated version holds at all time pairs for a continuous martingale
version. Killing only subtracts the left truncation energy. After killing,
$z=0$, so its dissipation also vanishes on the rest of a deterministic interval.

We carry out the fractional version of the probability iteration in
\cite[Sections 2--3]{HWW}. Fix $T>0$, $a\geq1$, and let
\[
 I_k=[(1-2^{-k})T,2T],\quad a_k=a(1-2^{-k}),\quad
 z_k=(\widehat\theta_n-a_k)_+,\quad
 U_k=\left(\int_{I_k}\norm{z_k(t)}_2^4\dd t\right)^{1/2}.
\]
For $k\geq1$, on $\{z_k>0\}$ one has
$z_{k-1}\geq a2^{-k}$ and $1+|\theta|\leq C2^kz_{k-1}$.
Let $X_{k-1}^*$ be the supremum of increments over time pairs in $I_{k-1}$
of the continuous martingale
\[
 X(t)=\sum_j\int_0^t\left(\int z_k g_j(x,\theta)\dd x\right)\dd W^j.
\]
Averaging the initial $L^2$ energy in
$I_{k-1}\setminus I_k$, whose length is $2^{-k}T$, and then integrating
\eqref{eq:truncation} gives
\begin{equation}\label{eq:iterationenergy}
 \sup_{I_k}\norm{z_k}_2^2+
 \int_{I_k}\norm{\Lambda^{1/2}z_k}_2^2\dd t
 \leq C_T^k(U_{k-1}+X_{k-1}^*).
\end{equation}
For this averaging, $\int_{I_{k-1}}\norm{z_{k-1}}_2^2\dd t
\leq(2T)^{1/2}U_{k-1}$. The trace drift is bounded by $C^k$ times that
integral, and every martingale difference is bounded by $X_{k-1}^*$.
Thus the selected initial time is chosen pathwise; it need not be a stopping
time.

The inhomogeneous periodic Sobolev inequality is
\begin{equation}\label{eq:fractionalSobolev}
 \norm{z}_4^2\leq C\big(\norm{\Lambda^{1/2}z}_2^2+\norm{z}_2^2\big).
\end{equation}
It follows from the homogeneous $H^{1/2}$ inequality after subtracting the
mean. The extra $L^2$ term is bounded in time by $(2T)^{1/2}U_{k-1}$.
Chebyshev's inequality gives
$|\{z_k(t)>0\}|\leq(2^k/a)^2\norm{z_{k-1}(t)}_2^2$.
Spatial H\"older, followed by time H\"older, therefore yields
\[
 U_k\leq(2^k/a)^{2/5}
       \norm{z_k}_{L^5(I_k;L^{5/2})}^2 U_{k-1}^{1/5}.
\]
The interpolation
$\norm{z}_{L^5_tL^{5/2}_x}
\leq\norm{z}_{L^\infty_tL^2_x}^{3/5}
     \norm{z}_{L^2_tL^4_x}^{2/5}$,
together with \eqref{eq:iterationenergy}--\eqref{eq:fractionalSobolev}, proves
\begin{equation}\label{eq:recurrence}
 U_k\leq C_T^k a^{-2\delta}
             (U_{k-1}+X_{k-1}^*)U_{k-1}^\delta,
 \qquad\delta=\frac15.
\end{equation}
This value uses the actual order-one fractional dissipation in dimension two.

By \eqref{eq:growth} and Cauchy--Schwarz in space, the bracket of $X$ over
$I_{k-1}$ is bounded by $C^k U_{k-1}^2$. Start this martingale at the
deterministic left endpoint of that interval. Lemma~\ref{lem:martingale} gives
\begin{equation}\label{eq:iterationprob}
 \PP\{X_{k-1}^*>\ell b,\ U_{k-1}\leq b\}
 \leq C\exp(-\ell^2/C^k),\qquad b,\ell>0.
\end{equation}
Only the implication from $U_{k-1}\leq b$ to the bracket bound is used.

Here is the threshold choice for \eqref{eq:recurrence}. Enlarge $B\geq2$
so the bracket constants are at most $B^k$, and set
\[
 b_k=a^2\lambda^k/\varrho^2,\qquad
 \ell_k=(2B)^{k/2}\varrho^\delta.
\]
Choose $\lambda\in(0,1)$ so that
$C_T(2B)^{1/2}\lambda^\delta\leq1/4$, and then choose
$\varrho\geq\varrho_0(T)$ sufficiently large. If
$U_{k-1}\leq b_{k-1}$ and $X_{k-1}^*\leq\ell_kb_{k-1}$, division of
\eqref{eq:recurrence} by $b_k$ bounds the result by
\[
 C_T^k(1+\ell_k)\lambda^{\delta(k-1)-1}\varrho^{-2\delta}
 \leq C\lambda^{-1-\delta}\varrho^{-\delta}\leq1.
\]
On $U_0\leq a^2/\varrho^2$, the first failed level must therefore lie in
an event in \eqref{eq:iterationprob}. Their probabilities sum to at most
$C\sum_{k\geq1}\exp(-c2^k\varrho^{2\delta})$. If none fails,
$U_k\to0$, which forces $(\widehat\theta_n-a)_+=0$ almost everywhere on
$[T,2T]\times\T^2$. Hence
\begin{equation}\label{eq:amplitudetail}
 \PP\left\{\norm{\widehat\theta_n^+}_{L^\infty([T,2T]\times\T^2)}>a,
                 \ \varrho U_0^{1/2}\leq a\right\}
 \leq C\exp(-c\varrho^\delta).
\end{equation}
All these constants are independent of $n$.

Take $\varrho=\sqrt a$ for sufficiently large $a$.
By \eqref{eq:L2moments}, $U_0$ has moments of every finite order, uniformly
in $n$. The probability in the complementary event is
$\PP\{U_0>a\}$, so \eqref{eq:amplitudetail} yields amplitude moments of
every finite order. Apply the same calculation to $-\theta$.
Continuity before killing turns the essential bounds into pointwise time and
space suprema. A finite covering by positive-time intervals proves
\eqref{eq:amplitudemoments}.

Monotone exhaustion of $\tau_n$ gives finite amplitude on
$[\varepsilon,S\wedge\tau_*)$ almost surely, simultaneously for rational
$\varepsilon>0$ and integer $S$. Since $\tau_*>0$, each sample admits an
integer $m$ with $1/m<\tau_*$. Local $H^3$ continuity bounds the solution
on $[0,1/m]$; the positive-time estimate bounds the remaining interval.
This proves \eqref{eq:amplitudeaslifetime}.
\end{proof}

\section{Spatial regularity and global continuation}\label{sec:regularity}

We now fix an integer amplitude level $M\geq1$. All estimates in this section
are initially stopped both at $\tau_n$ and at
$\sigma_M=\inf\{t<\tau_*:\norm{\theta(t)}_\infty\geq M\}$,
with the infimum of the empty set equal to infinity.
The exponents below are chosen deterministically for that $M$.

\subsection{Dissipation of finite powers of increments}

Write $u=R^\perp\theta$, $w(x,h)=\delta_h\theta(x)=\theta(x+h)-\theta(x)$,
using periodic extension in $h\in\R^2$. Then
\begin{equation}\label{eq:increment}
 \dd w+(u(x)\cdot\nabla_x+\delta_hu(x)\cdot\nabla_h+\Lambda_x)w\dd t
       =\sum_k\delta_hg_k\dd W^k,
\end{equation}
where $\delta_hg_k=g_k(x+h,\theta(x+h))-g_k(x,\theta(x))$.
Both transport fields are divergence free in their integration variables.
The bounded scalar and spatial derivatives imply
\begin{equation}\label{eq:noiseincrement}
 \left(\sum_k|\delta_hg_k|^2\right)^{1/2}\leq L_0|w|+B_0|h|.
\end{equation}
The constants are independent of derivatives of $\theta$.

The following elementary scalar inequality preserves the dependence on $p$
needed to choose a spatial exponent before selecting a large finite $p$.

\begin{lemma}\label{lem:convexity}
For real $a,b$ and $p\geq2$,
\begin{equation}\label{eq:convexity}
 |b|^p-|a|^p-p|a|^{p-2}a(b-a)
       \geq\frac p2|a|^{p-2}(a-b)^2.
\end{equation}
Consequently, with
\[
 D[w](x)=\frac1{2\pi}\int_{\R^2}
           \frac{|w(x)-w(x+y)|^2}{|y|^3}\dd y,
\]
one has
\begin{equation}\label{eq:pdissipation}
 p|w|^{p-2}w\Lambda w-\Lambda(|w|^p)
       \geq\frac p2|w|^{p-2}D[w].
\end{equation}
\end{lemma}

\begin{proof}
For $a\ne0$ divide by $|a|^p$ and put $t=b/a$. If $t\geq0$, then
\[
 t^p-1-p(t-1)-(p-1)(t-1)^2
 =t\big(t^{p-1}-(p-1)t+p-2\big)\geq0
\]
by convexity, and $p-1\geq p/2$. If $t=-s\leq0$, convexity of
$q\mapsto q^{p/2}$ gives $s^p\geq(p/2)s^2+1-p/2$ and hence
\[
 s^p+ps+p-1\geq(p/2)(s+1)^2.
\]
The case $a=0$ follows directly; at $p=2$ interpret $|a|^0=1$,
and the inequality is an equality.
Integrate \eqref{eq:convexity} against the positive kernel of $\Lambda$ to
obtain \eqref{eq:pdissipation}.
\end{proof}

We recall the spatial estimates of Constantin, Tarfulea and Vicol
\cite[(4.15)--(4.25), August 2, 2013 author version]{CTV}:
\begin{equation}\label{eq:kernelbounds}
 D[w]\geq c_0\frac{|w|^3}{M|h|},\qquad
 |\delta_hu|\leq C_0\left(\sqrt{rD[w]}+\frac{M|h|}{r}\right),
 \qquad r\geq4|h|.
\end{equation}
Their whole-$\R^2$ periodic-kernel proof also makes the role of the mean and
of large splitting radii explicit. For a radial cutoff $\eta$ equal to one
on $[0,1]$ and zero beyond $2$, put
$F_r(y)=(1-\eta(|y|/r))|y|^{-3}$. Expanding the square in $D$ and shifting
the finite difference in the cross term gives
\[
 D[w]\geq c|w|^2/r-C|w|M|h|/r^2,
 \qquad \norm{\delta_{-h}F_r}_1\leq C|h|/r^2.
\]
Selecting $r=c_{\rm r}M|h|/|w|$, with a sufficiently large universal
$c_{\rm r}$, proves the lower bound. For the velocity split the Riesz kernel
$y^\perp/(2\pi|y|^3)$ at radius $r$. The radial inner kernel has zero
integral. Subtracting $w(x)$ and applying Cauchy--Schwarz gives
$C\sqrt{rD[w]}$. In the outer part, shift $\delta_h$ to the kernel; its
derivative has $L^1$ norm at most $C/r$, giving $CM|h|/r$.
The far scalar kernel is integrable, and the shifted Riesz kernel is
integrable. With a large outer cutoff of radius $L$, the shift error is
$O(M|h|/L)$ and disappears as $L\to\infty$. These bounds hold for every
$r\geq4|h|$, including radii beyond a fundamental period. Constants cancel
in the increments, so no zero-mean condition is involved. Vanishing $w$ is
treated separately before division.

Take now $r=c_{\rm r}M|h|/|w|$ with a fixed $c_{\rm r}\geq8$. Since
$|w|\leq2M$, this radius is admissible. Young's inequality gives
\[
 \alpha\frac{|\delta_hu|}{|h|}|w|^2
 \leq\frac18D[w]+C\left(\alpha^2c_{\rm r}M+\frac\alpha{c_{\rm r}}\right)
                       \frac{|w|^3}{|h|}.
\]
The lower bound in \eqref{eq:kernelbounds} shows that a sufficiently small
universal $c_*>0$ allows
\begin{equation}\label{eq:absorption}
 0<\alpha<\min\{1/4,c_*/M\},\qquad
 \alpha\frac{|\delta_hu|}{|h|}|w|^2\leq\frac14D[w].
\end{equation}
The choice of $\alpha$ is independent of $p$. Multiplication by
$p|w|^{p-2}$ lets \eqref{eq:pdissipation} absorb the derivative of the
singular increment weight for any finite $p$.

\subsection{An integrated increment estimate}

Choose a nonnegative $\chi\in C_c^\infty(\R^2)$ equal to one near zero,
with support in a small ball. For a fixed even integer $p>8/\alpha$ define
\begin{equation}\label{eq:Y}
 Y(t)=\int_{\T^2}\int_{\R^2}
       \chi(h)\frac{|\delta_h\theta(x,t)|^p}{|h|^{\alpha p}}\dd h\dd x.
\end{equation}
The measure in $h$ is ordinary Lebesgue measure. Thus $Y$ is a truncated
fractional Sobolev seminorm of order $s=\alpha-2/p$, and the H\"older exponent
it controls in dimension two is
\begin{equation}\label{eq:gamma}
 \gamma=\alpha-4/p>0.
\end{equation}

\begin{proposition}\label{prop:holder}
For each fixed $M$ and the deterministic exponents in
\eqref{eq:absorption}--\eqref{eq:gamma}, the pre-lifetime $C^\gamma$ norm
is almost surely bounded on every finite horizon before the amplitude-$M$ stop.
The bounds obtained by local stopping are uniform in $n$.
\end{proposition}

\begin{proof}
Apply It\^o's formula to $|w|^p$ in \eqref{eq:increment} and integrate against
$\chi(h)|h|^{-\alpha p}$. Spatial transport integrates to zero, as does
$\int\Lambda_x(|w|^p)\dd x$. The weight's transport derivative contributes
\[
 -\alpha p\int\chi(h)|w|^p|h|^{-\alpha p-2}\delta_hu\cdot h\dd h\dd x,
\]
which is absorbed using \eqref{eq:pdissipation} and \eqref{eq:absorption}.
The derivative of $\chi$ is supported away from zero and has bounded
contribution depending on $M,p$, since $|w|\leq2M$ and
$\norm{\delta_hu}_1\leq C\norm{\theta}_2\leq CM$.

Let $J=\int\chi(h)|h|^{p(1-\alpha)}\dd h\dd x<\infty$.
By \eqref{eq:noiseincrement} and H\"older, the It\^o trace correction is
at most
\[
 C_p\left(Y+Y^{(p-2)/p}J^{2/p}\right)\leq C_p(1+Y),
\]
and the Euclidean norm of the integrated martingale coefficient is at most
\[
 pL_0Y+pB_0Y^{(p-1)/p}J^{1/p}\leq C_p(1+Y).
\]
We have therefore proved
\begin{equation}\label{eq:Yenergy}
 \dd Y\leq C_{M,p}(1+Y)\dd t+\dd\mathcal N,
 \qquad\dd[\mathcal N]_t\leq C_p(1+Y)^2\dd t.
\end{equation}

For completeness, all integrations are first carried out with $|h|>\epsilon$
at a fixed local $H^3$ stop. Then $w=O(|h|)$ and $\delta_hu=O(|h|)$,
so the inner boundary flux is $O(\epsilon^{2+p(1-\alpha)})$.
Choose $1/2<\beta<1$. Local $C^{1,\beta}$ regularity gives
\[
 |\delta_yw|\leq C_n|h|\min\{|y|^\beta,1\},\qquad D[w]\leq C_n|h|^2.
\]
Here the near-$y$ integral is bounded by
$\int_0^1 r^{2\beta-2}\dd r<\infty$. After spatial cancellation, the
remaining drift and dissipation have integrable bounds of order
$|h|^{p(1-\alpha)}$; the noise terms have the same integrable powers.
Dominated convergence and the stopped stochastic isometry remove $\epsilon$
before any uniform limit in $n$ is taken. Thus \eqref{eq:Yenergy} holds for
$Y$ itself at the local and amplitude stops.

The datum is smooth, so $Y(0)<\infty$. Lemma~\ref{lem:martingale} supplies
a tail tending to zero for its stopped supremum uniformly in $n$. The
suprema on the increasing stopped intervals are increasing, so their limit
is finite almost surely.

We spell out the embedding to keep track of the loss in \eqref{eq:gamma}.
For a sufficiently small torus ball $B_r$, all pairwise differences lie where
$\chi=1$. Jensen's inequality and \eqref{eq:Y} give
\[
 \fint_{B_r}|\theta-\theta_{B_r}|^p
 \leq |B_r|^{-2}\int_{B_r}\int_{B_r}|\theta(x)-\theta(y)|^p\dd x\dd y
 \leq C r^{\alpha p-4}Y=C r^{p\gamma}Y.
\]
Differences of averages on concentric radii $r,r/2$ are bounded by
$Cr^\gamma Y^{1/p}$. Sum the geometric series and use local continuity to
obtain $|\theta(x)-\theta_{B_r(x)}|\leq Cr^\gamma Y^{1/p}$.
For $|x-y|=r$, compare these averages in a common ball of radius $3r$.
For larger distances use the amplitude bound. Consequently
\begin{equation}\label{eq:holderembedding}
 \norm{\theta}_{C^\gamma}\leq C_{\alpha,p,\chi}(M+Y^{1/p}),
\end{equation}
which proves the proposition.
\end{proof}

\subsection{The high-regularity estimate}

We need a continuation bound that uses only a positive H\"older exponent and
the amplitude. Fix a periodic inhomogeneous Littlewood--Paley decomposition
$(\Delta_j)_{j\geq-1}$ and put
$E_s=\sum_{j\geq-1}2^{2sj}\norm{\Delta_j\theta}_2^2\asymp\norm{\theta}_{H^s}^2$.
The low block includes the constant mode.

\begin{lemma}\label{lem:transport}
For smooth real periodic $\theta$, $u=R^\perp\theta$, $0<\gamma<1$
and $s\geq3$,
\begin{equation}\label{eq:dyadictransport}
 \left|\sum_j2^{2sj}
  \langle\Delta_j(u\cdot\nabla\theta),\Delta_j\theta\rangle\right|
 \leq C_{s,\gamma}\norm{\theta}_{C^\gamma}
                    \norm{\theta}_{H^{s+(1-\gamma)/2}}^2.
\end{equation}
For $\norm{\theta}_\infty\leq M$, one also has
\begin{equation}\label{eq:tameH3}
 \sum_k\norm{g_k(\cdot,\theta)}_{H^3}^2
       \leq C_M(1+\norm{\theta}_{H^3}^2).
\end{equation}
\end{lemma}

\begin{proof}
Let $\mathcal K=\norm{\theta}_{C^\gamma}$, and write $\theta_j=\Delta_j\theta$
and $S_{j-1}=\sum_{\ell\leq j-2}\Delta_\ell$. The periodic block kernels
are periodizations of Schwartz kernels; their $L^1$ norms are bounded and
their first moments are $O(2^{-j})$. Their zero average, also for the smooth
annular multiplier $R\Delta_j$, gives
\[
 \norm{\theta_j}_\infty+\norm{u_j}_\infty
       \leq C\mathcal K2^{-\gamma j},\qquad
 \norm{\nabla S_{j-1}u}_\infty+
 \norm{\nabla S_{j-1}\theta}_\infty
       \leq C_\gamma\mathcal K2^{(1-\gamma)j}.
\]
The finitely many low modes satisfy the corresponding bound directly.

Decompose the product into low--high, high--low and neighboring high--high
terms. In low--high interactions $|k-j|\leq C$. Commuting $\Delta_j$
past $S_{k-1}u\cdot\nabla$ costs
\[
 C2^{-j}\norm{\nabla S_{k-1}u}_\infty\,2^k\norm{\theta_k}_2.
\]
The remaining principal transport is $S_{j-1}u\cdot\nabla\theta_j$,
whose energy pairing is zero. Its difference from $S_{k-1}u$ costs
$C\mathcal K2^{-\gamma j}$ times the derivative of the neighboring block.
High--low terms cost
$\norm{u_k}_2\norm{\nabla S_{k-1}\theta}_\infty$ for the same neighboring
indices. These terms are bounded by
$C\mathcal K\norm{\theta}_{H^{s+(1-\gamma)/2}}^2$.

For high--high terms use $u\cdot\nabla\theta=\nabla\cdot(u\theta)$.
The block norm is bounded by
\[
 C\mathcal K2^j\sum_{k\geq j-C}2^{-\gamma k}\norm{\theta_k}_2.
\]
Set $\sigma=s+(1-\gamma)/2$ and $b_j=2^{\sigma j}\norm{\theta_j}_2$.
The weighted pairings sum to at most
\[
 C\mathcal K\sum_{j\geq-1}\sum_{k\geq j-C}
       2^{-(s+(1+\gamma)/2)(k-j)}b_jb_k
       \leq C\mathcal K\sum_jb_j^2.
\]
The convolution is summable. The low block costs at most $C\mathcal K E_s$.
This proves \eqref{eq:dyadictransport}; the proof also bounds the sum of
absolute block pairings after the indicated cancellations.

For \eqref{eq:tameH3}, the third-derivative chain rule contains only the
nonlinear highest-order products $D\theta D^2\theta$ and $(D\theta)^3$.
Let $\mathfrak a=\norm{D\theta}_6$, $\mathfrak b=\norm{D^2\theta}_3$
and $\mathfrak c=\norm{D^3\theta}_2$. Periodic integration by parts gives
\[
 \mathfrak a^6\leq CM\mathfrak a^4\mathfrak b,
 \qquad \mathfrak b^3\leq C\mathfrak a\mathfrak b\mathfrak c.
\]
For the first inequality integrate
$\operatorname{div}(|D\theta|^4D\theta)$ against $\theta$.
For the second integrate one derivative in
$\sum_{i,j}\int\theta_{ij}^2|D^2\theta|$, and use
$1/6+1/3+1/2=1$; a smooth regularization of the Hessian length justifies
its value at zero. Unless one of the derivative norms vanishes, in which
case the estimates are immediate, these inequalities imply
\[
 \mathfrak a\leq CM^{2/3}\mathfrak c^{1/3},\qquad
 \mathfrak b\leq CM^{1/3}\mathfrak c^{2/3}.
\]
Thus $\norm{D\theta D^2\theta}_2\leq\mathfrak a\mathfrak b\leq C_M\mathfrak c$
and $\norm{(D\theta)^3}_2=\mathfrak a^3\leq C_M\mathfrak c$.
All mixed spatial coefficient derivatives are bounded by assumption; lower
products obey the same bounds and finite-volume H\"older. Finally
$\norm{g_k(\cdot,\theta)}_2\leq C(1+M)$, proving \eqref{eq:tameH3}.
\end{proof}

Apply It\^o's formula to $E_3$, first summing only finitely many whole dyadic
blocks. Each such block has finite Fourier support on the torus. At a local
$H^3$ stop, the nonlinear pairings converge absolutely in time by the preceding
proof and $\theta\in L^2_tH^{7/2}$; the dissipation converges monotonically.
The trace and martingale terms converge by \eqref{eq:tameH3} and the stopped
isometry. This proves the energy formula for $E_3$ without assuming an
$H^3$-valued drift.

The dissipation bounds the high-frequency $H^{7/2}$ norm. The finitely many
low modes cost $CE_3$, including the nondissipated constant mode.
Interpolation and Young's inequality give
\[
 C\mathcal K\norm{\theta}_{H^{3+(1-\gamma)/2}}^2
 \leq\varepsilon\norm{\theta}_{H^{7/2}}^2
       +C_{\gamma,\varepsilon}\mathcal K^{1/\gamma}\norm{\theta}_{H^3}^2.
\]
Consequently, before the amplitude-$M$ and $C^\gamma$ level-$R$ stops,
\begin{equation}\label{eq:H3energy}
 \dd E_3\leq C_{M,R,\gamma}(1+E_3)\dd t+\dd\mathcal M,
 \qquad \dd[\mathcal M]_t\leq C_M(1+E_3)^2\dd t.
\end{equation}
For the bracket use Cauchy--Schwarz in the dyadic Hilbert norm and
\eqref{eq:tameH3}. Lemma~\ref{lem:martingale} now bounds the stopped
$\sup E_3$ by tails independent of the local level $n$.

\subsection{Exclusion of a finite lifetime}

\begin{proof}[Proof of Theorem~\ref{thm:main}]
For each integer $M\geq1$ fix deterministic $\alpha_M$ satisfying
\eqref{eq:absorption}, and then an even integer $p_M>8/\alpha_M$.
Put $\gamma_M=\alpha_M-4/p_M>0$. Proposition~\ref{prop:holder} yields a
finite pre-lifetime $C^{\gamma_M}$ supremum at the amplitude-$M$ stop on
every finite horizon. The uniform local-stop tail and monotone exhaustion
justify this assertion. Take a common probability-one event over integers
$M$, finite integer horizons, and the later H\"older levels $R$ in
\eqref{eq:H3energy}.

On a sample with finite $\tau_*$, Proposition~\ref{prop:amplitude} supplies
a finite amplitude supremum up to that lifetime. Some integer $M$ strictly
exceeds it. For this deterministic family member, the H\"older supremum is
finite, so some integer $R$ strictly exceeds it. Neither of these two stops
is then reached before $\tau_*$. The uniform $H^3$ tail in
\eqref{eq:H3energy}, exhausted monotonically in $n$, gives a finite
pre-lifetime $H^3$ supremum. This contradicts Proposition~\ref{prop:local}.
Hence $\PP\{\tau_*<\infty\}=0$.

Every exponent was fixed before its It\^o estimate. The sample-dependent
choice of one integer $M$, and then $R$, occurs only after the countable
intersection and introduces no anticipating stochastic integrand.

On every finite interval the patched local solution is continuous in $H^3$
and adapted to the original basis. Local pathwise uniqueness gives global
pathwise uniqueness by increasing common norm stops. Its drift is continuous
in $H^2$, since $H^2(\T^2)$ is an algebra. Its noise is predictable and
locally square integrable in $H^3$, so its stochastic integral has a continuous
$H^3$ version \cite[Proposition 2.1]{BVNVW}. The local weak identity holds
on a countable Fourier basis and rational times. Both sides are continuous
in $H^2$, so Fourier completeness and time continuity give the integral
identity at every time in $H^2$, hence in $H^1$. This proves all the assertions
of the theorem.
\end{proof}


\begin{thebibliography}{9}
\bibitem{AV}
A. Agresti and M. Veraar,
\emph{Nonlinear SPDEs and maximal regularity: an extended survey},
NoDEA Nonlinear Differential Equations Appl. \textbf{32} (2025), article 123.
\href{https://doi.org/10.1007/s00030-025-01090-2}{doi:10.1007/s00030-025-01090-2}.

\bibitem{BVNVW}
Z. Brze\'zniak, J. M. A. M. van Neerven, M. C. Veraar and L. Weis,
\emph{It\^o's formula in UMD Banach spaces and regularity of solutions of the Zakai equation},
J. Differential Equations \textbf{245} (2008), 30--58.
\href{https://doi.org/10.1016/j.jde.2008.03.026}{doi:10.1016/j.jde.2008.03.026}.
\href{https://arxiv.org/abs/0804.0302v1}{arXiv:0804.0302v1}.

\bibitem{CZ}
Z.-Q. Chen and X. Zhang,
\emph{Harnack inequalities for nonlocal operators with supercritical drifts and their applications},
preprint (2025), \href{https://arxiv.org/abs/2511.13012v1}{arXiv:2511.13012v1}.

\bibitem{CTV}
P. Constantin, A. Tarfulea and V. Vicol,
\emph{Long time dynamics of forced critical SQG},
Comm. Math. Phys. \textbf{335} (2015), 93--141.
\href{https://doi.org/10.1007/s00220-014-2129-3}{doi:10.1007/s00220-014-2129-3}.
The equation numbers cited here refer to the
\href{https://cims.nyu.edu/~vicol/CTV2.pdf}{August 2, 2013 author preprint}.

\bibitem{HWW}
E. P. Hsu, Y. Wang and Z. Wang,
\emph{Stochastic De Giorgi iteration and regularity of stochastic partial differential equations},
Ann. Probab. \textbf{45} (2017), no. 5, 2855--2866.
\href{https://doi.org/10.1214/16-AOP1126}{doi:10.1214/16-AOP1126}.
\href{https://arxiv.org/abs/1312.3311v3}{arXiv:1312.3311v3}.

\bibitem{RZZ}
M. R\"ockner, R. Zhu and X. Zhu,
\emph{Local existence and non-explosion of solutions for stochastic fractional partial differential equations driven by multiplicative noise},
Stochastic Process. Appl. \textbf{124} (2014), 1974--2002.
\href{https://doi.org/10.1016/j.spa.2014.01.010}{doi:10.1016/j.spa.2014.01.010}.
The source version used for theorem numbering is
\href{https://arxiv.org/abs/1307.4392v1}{arXiv:1307.4392v1}.

\bibitem{AIM478}
AIM problem catalogue,
\emph{478. Global regularity for critical SQG with state-dependent noise},
\href{https://github.com/MColbrook/AIM/blob/a153f469ee0ea64e876b2ad97c79d252c19a3e22/problems/478-critical-stochastic-sqg-global-regularity.md}{immutable problem statement, revision \texttt{a153f469ee0e}}.
\end{thebibliography}

\section*{AI use}
OpenAI Codex was used in the literature search, development and checking of
the mathematical argument, and preparation of this manuscript. Fresh-context
AI reader and referee reviews were used for exposition and source checking.
The author is responsible for the mathematical claims and citations.

\end{document}
